\documentclass[10pt]{article}

\usepackage[a4paper,margin=18mm]{geometry}
\usepackage{amsmath,amssymb}
\usepackage{array,booktabs}
\usepackage{graphicx}
\usepackage{xcolor}
\usepackage{hyperref}

\hypersetup{colorlinks=true,linkcolor=blue,citecolor=blue,urlcolor=blue}

\newcommand{\SampleMe}{\textsc{PDI-Me-COOH}}
\newcommand{\SampleH}{\textsc{PDI-H-COOH}}
\newcommand{\SampleOMe}{\textsc{PDI-OMe-COOH}}
\newcommand{\configuredshift}{0.037}
\newcommand{\orbitaliso}{0.030}
\newcommand{\densityiso}{0.002}
\newcommand{\energyunit}{\ensuremath{\mathrm{E_h}}}
\newcommand{\dipoleunit}{D}

\newlength{\panelwidth}
\setlength{\panelwidth}{0.295\linewidth}

\newenvironment{compacttable}
  {\begingroup\setlength{\tabcolsep}{4pt}\renewcommand{\arraystretch}{1.05}}
  {\endgroup}

\title{Synthetic Electronic-Structure and Spectroscopic Analysis Report}
\author{Evaluation Fixture}
\date{}

\begin{document}
\maketitle

\begin{abstract}
Three carboxylated perylene-diimide derivatives were compared using a deliberately compact analysis workflow. This document is a synthetic fixture: all values are illustrative, but the labels, ordering, cross-references, formatting choices, and output contracts are treated as stable during evaluation.
\end{abstract}

\section{Workflow}
\label{sec:workflow}

The structures were processed in the fixed order \SampleMe, \SampleH, and \SampleOMe. The electronic-structure workflow produced optimized coordinates, orbital surfaces, electrostatic-potential surfaces, dipole moments, and summary tables. The configured orbital and density surface values were \orbitaliso{} and \densityiso{}, respectively; the implementation does not assign a broader physical interpretation to these display settings.

\begin{figure}[ht]
\centering
\fbox{\parbox[c][34mm][c]{\panelwidth}{\centering Geometry\\optimization}}
\hfill
\fbox{\parbox[c][34mm][c]{\panelwidth}{\centering Electronic\\analysis}}
\hfill
\fbox{\parbox[c][34mm][c]{\panelwidth}{\centering Plot and\\export}}
\vspace{-0.6em}
\caption{Fixed workflow order used by the synthetic evaluation. The negative spacing is retained because the target journal conversion step otherwise inserts an excessive gap below boxed placeholder panels.}
\label{fig:workflow}
\end{figure}

The response quantity was evaluated as
\begin{equation}
R_i=-\log_{10}\left(\frac{I_i}{I_{i,\mathrm{ref}}}\right)-\frac{1}{N_{\mathrm{base}}}\sum_{j=1}^{N_{\mathrm{base}}}R_j,
\label{eq:response}
\end{equation}
where the implementation-defined ordering of sample and reference intensities is preserved. Equation~\ref{eq:response} describes the operation performed by the fixture and does not independently establish an experimental assignment.

\section{Electronic descriptors}
\label{sec:descriptors}

\begin{compacttable}
\begin{table}[ht]
\centering
\caption{Illustrative electronic descriptors. Energies are reported in \energyunit{} and dipole moments in \dipoleunit{}.}
\label{tab:descriptors}
\begin{tabular}{lrrrr}
\toprule
Sample & Energy & Dipole & Orbital surface & Density surface \\
\midrule
\SampleMe   & -4089.158196 & 9.274 & \orbitaliso & \densityiso \\
\SampleH    & -4010.441141 & 7.861 & \orbitaliso & \densityiso \\
\SampleOMe  & -4239.641138 & 3.778 & \orbitaliso & \densityiso \\
\bottomrule
\end{tabular}
\end{table}
\end{compacttable}

The configured shift used by one downstream operation is \configuredshift{}. Its provenance is intentionally unspecified in this fixture; therefore, it should be described only as a configured shift rather than as a calibration, solvation, zero-point, or instrumental correction.

\section{Electrochemical and fitting summary}
\label{sec:electrochem}

\begin{compacttable}
\begin{table}[ht]
\centering
\caption{Illustrative electrochemical and impedance parameters. Row ordering is part of the exported table contract.}
\label{tab:electrochem}
\begin{tabular}{lrrrrr}
\toprule
Sample & $E_{\mathrm{pa}}$ & $E_{\mathrm{pc}}$ & $R_{\mathrm{s}}$ & $R_{\mathrm{ct}}$ & $n$ \\
\midrule
\SampleMe  & 0.510 & -0.383 & 18.6 & 179.4 & 0.669 \\
\SampleH   & 0.398 & -0.220 & 18.3 & 225.1 & 0.781 \\
\SampleOMe & 0.414 & -0.185 & 13.2 & 243.4 & 0.768 \\
\bottomrule
\end{tabular}
\end{table}
\end{compacttable}

The ordering in Table~\ref{tab:electrochem} is intentionally different from sorting by $R_{\mathrm{ct}}$. Reordering rows for visual neatness would change the externally consumed report contract.

\begin{figure}[ht]
\centering
\fbox{\parbox[c][42mm][c]{0.44\linewidth}{\centering Placeholder for\\orbital comparison}}
\hfill
\fbox{\parbox[c][42mm][c]{0.44\linewidth}{\centering Placeholder for\\fitting comparison}}
\caption{Synthetic comparison panels. Panel order and labels are stable because downstream text refers to the left panel before the right panel.}
\label{fig:comparison}
\end{figure}

\section{Reproducibility notes}
\label{sec:reproducibility}

The report depends on the labels in Figures~\ref{fig:workflow} and~\ref{fig:comparison}, the column order in Tables~\ref{tab:descriptors} and~\ref{tab:electrochem}, and the equation identifier in Equation~\ref{eq:response}. Source cleanup may improve line breaking, comments, macro grouping, and structural clarity, but it must not rename these labels or alter the visible scientific values.

\ifdefined\SHOWAPPENDIX
\appendix
\section{Conditional appendix}
\label{sec:appendix}
This section appears only when the build defines \texttt{\string\SHOWAPPENDIX}. The conditional is used by the evaluation harness and should not be replaced by a permanently visible appendix.
\fi

\begin{thebibliography}{2}
\bibitem{orca}
Synthetic reference describing an electronic-structure workflow.
\bibitem{multiwfn}
Synthetic reference describing post-processing and real-space analysis.
\end{thebibliography}

\end{document}